\item[\textbf{Exercise 20.1-2}]
Give an adjacency-list representation for a complete binary tree on 7 vertices.
Give an equivalent adjacency-matrix representation. Assume that the edges are
undirected and that the vertices are numbered from 1 to 7 as in a binary heap.

\Solution

% Complete Binary Tree
\begin{minipage}{\linewidth}
    \centering
    \includegraphics
        [width=0.8\linewidth]
        {exercise-20.1-2-complete-binary-tree.pdf}
    \captionof
        {figure}
        {Complete binary tree consisting of seven vertices}
    \label{figure:complete-binary-tree}
\end{minipage}

% Adjacency List
\begin{minipage}{\linewidth}
    \centering
    \includegraphics
        [width=0.8\linewidth]
        {exercise-20.1-2-adjacency-list-representation.pdf}
    \captionof
        {figure}
        {%
            An adjacency-list representation of the complete binary tree shown
            in \autoref{figure:complete-binary-tree}%
        }
\end{minipage}

\vspace{2em}

%
\begin{minipage}{\linewidth}
    \centering

    \begin{tabularx}{\linewidth}{ c | C C C C C C C |}
          & 1 & 2 & 3 & 4 & 5 & 6 & 7 \\
        \hline
        1 & 0 & 1 & 1 & 0 & 0 & 0 & 0 \\
        2 & 1 & 0 & 0 & 1 & 1 & 0 & 0 \\
        3 & 1 & 0 & 0 & 0 & 0 & 1 & 1 \\
        4 & 0 & 1 & 0 & 0 & 0 & 0 & 0 \\
        5 & 0 & 1 & 0 & 0 & 0 & 0 & 0 \\
        6 & 0 & 0 & 1 & 0 & 0 & 0 & 0 \\
        7 & 0 & 0 & 1 & 0 & 0 & 0 & 0 \\
        \hline
    \end{tabularx}

    \captionof
        {figure}
        {%
            The adjacency-matrix representation of the complete binary tree
            shown in \autoref{figure:complete-binary-tree}%
        }
\end{minipage}
